\documentclass[11pt,a4paper]{article}

\usepackage[utf8]{inputenc}
\usepackage[T1]{fontenc}
\usepackage[french]{babel}
\usepackage[a4paper,margin=2.3cm]{geometry}
\usepackage{amsmath,amssymb,amsthm,mathtools}
\usepackage{lmodern}
\usepackage{enumitem}
\usepackage{hyperref}
\usepackage{xcolor}
\usepackage{fancyhdr}
\usepackage{stmaryrd}
\usepackage{mathrsfs}

\hypersetup{
    colorlinks=true,
    linkcolor=black,
    urlcolor=blue,
    citecolor=black
}

\setlength{\parindent}{0pt}
\setlength{\parskip}{0.6em}

\pagestyle{fancy}
\fancyhf{}
\fancyhead[L]{Algèbre linéaire --- Chapitre 5}
\fancyhead[R]{Espaces vectoriels}
\fancyfoot[C]{\thepage}

% Environnements
\newtheorem{definition}{Définition}[section]
\newtheorem{theoreme}[definition]{Théorème}
\newtheorem{proposition}[definition]{Proposition}
\newtheorem{corollaire}[definition]{Corollaire}
\newtheorem{lemme}[definition]{Lemme}
\theoremstyle{remark}
\newtheorem{remarque}[definition]{Remarque}
\newtheorem{exemple}[definition]{Exemple}
\newtheorem*{methode}{Méthode}
\newtheorem*{idee}{Idée}

% Commandes
\newcommand{\R}{\mathbb{R}}
\newcommand{\N}{\mathbb{N}}
\newcommand{\Z}{\mathbb{Z}}
\newcommand{\Q}{\mathbb{Q}}
\newcommand{\C}{\mathbb{C}}
\newcommand{\K}{\mathbb{K}}
\newcommand{\PP}{\mathcal{P}}

\begin{document}

\begin{center}
    {\LARGE \textbf{Chapitre 5 --- Espaces vectoriels}}\\[0.5em]
    {\large Objet central de l’algèbre linéaire}
\end{center}

\hrule
\vspace{1em}

\tableofcontents

\newpage

\section{Pourquoi cette notion est-elle fondamentale ?}

L’idée d’un espace vectoriel est de regrouper dans un même cadre des objets très différents :
\begin{itemize}
    \item des vecteurs du plan ou de l’espace ;
    \item des \(n\)-uplets de réels ;
    \item des polynômes ;
    \item des fonctions ;
    \item des matrices.
\end{itemize}

Tous ces objets ont un point commun : on peut
\begin{itemize}
    \item les additionner ;
    \item les multiplier par un scalaire ;
    \item effectuer sur eux des combinaisons linéaires.
\end{itemize}

L’espace vectoriel est donc la structure abstraite qui permet d’unifier une grande partie des mathématiques du programme.

\begin{remarque}
Dans tout ce chapitre, le corps des scalaires sera en général \(\R\) ou \(\C\).  
Plus généralement, on pourra travailler sur un corps \(\K\), souvent égal à \(\R\) ou \(\C\).
\end{remarque}

\section{Définition d’un espace vectoriel}

\subsection{Définition générale}

\begin{definition}
Soit \(\K\) un corps.  
Un \textbf{espace vectoriel sur \(\K\)} est un ensemble \(E\) muni :
\begin{itemize}
    \item d’une loi de composition interne, appelée \textbf{addition},
    \[
    + : E\times E \to E,\qquad (u,v)\mapsto u+v,
    \]
    \item d’une loi de composition externe, appelée \textbf{multiplication par un scalaire},
    \[
    \K\times E \to E,\qquad (\lambda,u)\mapsto \lambda u,
    \]
\end{itemize}
telles que, pour tous \(u,v,w\in E\) et tous \(\lambda,\mu\in\K\), les propriétés suivantes soient vérifiées :
\begin{enumerate}[label=\textbf{(EV\arabic*)}]
    \item \(u+v=v+u\) ;
    \item \((u+v)+w=u+(v+w)\) ;
    \item il existe un élément \(0_E\in E\) tel que, pour tout \(u\in E\),
    \[
    u+0_E=u ;
    \]
    \item pour tout \(u\in E\), il existe un élément \(-u\in E\) tel que
    \[
    u+(-u)=0_E ;
    \]
    \item \(1u=u\) ;
    \item \((\lambda\mu)u=\lambda(\mu u)\) ;
    \item \(\lambda(u+v)=\lambda u+\lambda v\) ;
    \item \((\lambda+\mu)u=\lambda u+\mu u\).
\end{enumerate}
Les éléments de \(E\) sont appelés \textbf{vecteurs}, et les éléments de \(\K\) sont appelés \textbf{scalaires}.
\end{definition}

\begin{remarque}
Cette définition peut sembler longue, mais elle formalise simplement les règles naturelles du calcul vectoriel.
\end{remarque}

\subsection{Sens concret de la définition}

\begin{idee}
Un espace vectoriel est un ensemble dans lequel :
\begin{itemize}
    \item on peut additionner deux objets du même type ;
    \item on peut les dilater, les contracter, les annuler, les changer de sens par multiplication par un scalaire ;
    \item ces opérations se comportent de façon cohérente.
\end{itemize}
\end{idee}

\begin{exemple}
Dans le plan usuel, deux vecteurs peuvent être additionnés, et un vecteur peut être multiplié par un réel. On obtient encore un vecteur du plan. Le plan vectoriel est donc un espace vectoriel réel.
\end{exemple}

\section{Exemples fondamentaux}

\subsection{L’espace \(\K^n\)}

\begin{proposition}
Pour tout entier \(n\geq 1\), l’ensemble \(\K^n\) muni de l’addition terme à terme et de la multiplication par un scalaire terme à terme est un espace vectoriel sur \(\K\).
\end{proposition}

\begin{proof}
Soient
\[
u=(x_1,\dots,x_n),\qquad v=(y_1,\dots,y_n)\in \K^n,
\]
et \(\lambda\in\K\). On définit
\[
u+v=(x_1+y_1,\dots,x_n+y_n),
\qquad
\lambda u=(\lambda x_1,\dots,\lambda x_n).
\]
Les propriétés d’espace vectoriel découlent directement des propriétés usuelles des opérations dans \(\K\), appliquées coordonnée par coordonnée.
\end{proof}

\begin{exemple}
\[
\R^2,\qquad \R^3,\qquad \C^n
\]
sont des espaces vectoriels.
\end{exemple}

\subsection{L’ensemble des matrices}

\begin{proposition}
L’ensemble \(M_{n,p}(\K)\) des matrices à \(n\) lignes et \(p\) colonnes à coefficients dans \(\K\) est un espace vectoriel sur \(\K\).
\end{proposition}

\begin{proof}
L’addition des matrices et la multiplication d’une matrice par un scalaire se font coefficient par coefficient. Les axiomes d’espace vectoriel en découlent immédiatement.
\end{proof}

\begin{exemple}
\[
M_2(\R),\qquad M_{3,1}(\R),\qquad M_{n,p}(\C)
\]
sont des espaces vectoriels.
\end{exemple}

\subsection{L’ensemble des polynômes}

\begin{proposition}
L’ensemble \(\K[X]\) des polynômes à coefficients dans \(\K\) est un espace vectoriel sur \(\K\).
\end{proposition}

\begin{proof}
La somme de deux polynômes est un polynôme, et le produit d’un polynôme par un scalaire est encore un polynôme. Les axiomes proviennent des propriétés des coefficients dans \(\K\).
\end{proof}

\begin{exemple}
Si
\[
P(X)=X^2+2X+1,\qquad Q(X)=3X^3-X,
\]
alors
\[
P+Q=3X^3+X^2+X+1,
\qquad
2P=2X^2+4X+2.
\]
\end{exemple}

\subsection{L’ensemble des fonctions}

\begin{proposition}
Soit \(A\) un ensemble. L’ensemble \(\mathcal{F}(A,\K)\) des fonctions de \(A\) dans \(\K\) est un espace vectoriel sur \(\K\), muni des opérations
\[
(f+g)(x)=f(x)+g(x),\qquad (\lambda f)(x)=\lambda f(x).
\]
\end{proposition}

\begin{proof}
La somme de deux fonctions de \(A\) dans \(\K\) est encore une fonction de \(A\) dans \(\K\), de même que le produit d’une fonction par un scalaire. Les axiomes se vérifient point par point.
\end{proof}

\begin{exemple}
L’ensemble des fonctions réelles sur \(\R\), noté \(\mathcal{F}(\R,\R)\), est un espace vectoriel réel.

L’ensemble des fonctions continues sur un intervalle \(I\), noté \(\mathcal{C}(I,\R)\), est aussi un espace vectoriel réel.
\end{exemple}

\subsection{Exemples à ne pas confondre}

\begin{exemple}
L’ensemble
\[
\R_+=\{x\in\R\mid x\geq 0\}
\]
n’est pas un espace vectoriel sur \(\R\), car si \(x>0\), alors
\[
(-1)x=-x<0,
\]
donc \(-x\notin \R_+\).
\end{exemple}

\begin{exemple}
L’ensemble des fonctions positives sur \(\R\) n’est pas un espace vectoriel réel, car le produit par \(-1\) fait sortir de cet ensemble.
\end{exemple}

\section{Premières propriétés générales}

\subsection{Unicité du vecteur nul}

\begin{proposition}
Dans un espace vectoriel \(E\), le vecteur nul est unique.
\end{proposition}

\begin{proof}
Supposons qu’il existe deux vecteurs nuls \(0_E\) et \(0'_E\).  
Comme \(0_E\) est neutre pour l’addition,
\[
0_E+0'_E=0'_E.
\]
Comme \(0'_E\) est aussi neutre,
\[
0_E+0'_E=0_E.
\]
Donc
\[
0_E=0'_E.
\]
\end{proof}

\subsection{Unicité de l’opposé}

\begin{proposition}
Tout vecteur \(u\in E\) admet un unique opposé.
\end{proposition}

\begin{proof}
Soient \(v\) et \(w\) deux opposés de \(u\). Alors
\[
u+v=0_E
\qquad \text{et} \qquad
u+w=0_E.
\]
On a
\[
v=v+0_E=v+(u+w)=(v+u)+w=(u+v)+w=0_E+w=w.
\]
Donc l’opposé de \(u\) est unique.
\end{proof}

\subsection{Calculs usuels dans un espace vectoriel}

\begin{proposition}
Pour tout vecteur \(u\in E\) et tout scalaire \(\lambda\in\K\), on a :
\begin{enumerate}
    \item \(0u=0_E\) ;
    \item \(\lambda 0_E=0_E\) ;
    \item \((-1)u=-u\) ;
    \item si \(\lambda u=0_E\), alors \(\lambda=0\) ou \(u=0_E\).
\end{enumerate}
\end{proposition}

\begin{proof}
\textbf{1.} On a
\[
(0+0)u=0u+0u.
\]
Or \(0+0=0\), donc
\[
0u=0u+0u.
\]
En ajoutant l’opposé de \(0u\) aux deux membres, on obtient
\[
0_E=0u.
\]

\textbf{2.} On a
\[
\lambda(0_E+0_E)=\lambda 0_E+\lambda 0_E.
\]
Comme \(0_E+0_E=0_E\),
\[
\lambda 0_E=\lambda 0_E+\lambda 0_E.
\]
En ajoutant l’opposé de \(\lambda 0_E\), on obtient
\[
0_E=\lambda 0_E.
\]

\textbf{3.} On a
\[
u+(-1)u=(1+(-1))u=0u=0_E.
\]
Ainsi \((-1)u\) est l’opposé de \(u\), donc
\[
(-1)u=-u.
\]

\textbf{4.} Supposons \(\lambda u=0_E\).  
Si \(\lambda\neq 0\), alors \(\lambda\) admet un inverse \(\lambda^{-1}\) dans \(\K\). En multipliant par \(\lambda^{-1}\), on obtient
\[
u=\lambda^{-1}(\lambda u)=\lambda^{-1}0_E=0_E.
\]
Donc \(\lambda=0\) ou \(u=0_E\).
\end{proof}

\begin{remarque}
Le dernier point ressemble à la propriété de produit nul pour les scalaires, mais il faut être prudent : ici, on l’obtient grâce à l’existence de l’inverse de \(\lambda\) quand \(\lambda\neq 0\).
\end{remarque}

\section{Sous-espaces vectoriels}

\subsection{Définition}

\begin{definition}
Soit \(E\) un espace vectoriel sur \(\K\).  
Un \textbf{sous-espace vectoriel} de \(E\) est une partie \(F\subset E\) qui est elle-même un espace vectoriel pour les lois induites par \(E\).
\end{definition}

\begin{remarque}
En pratique, on ne vérifie presque jamais directement tous les axiomes. Il existe un critère simple.
\end{remarque}

\subsection{Critère de sous-espace vectoriel}

\begin{theoreme}[Critère de sous-espace vectoriel]
Soit \(F\subset E\). Alors \(F\) est un sous-espace vectoriel de \(E\) si et seulement si :
\begin{enumerate}
    \item \(F\neq \varnothing\) ;
    \item pour tous \(u,v\in F\), on a \(u+v\in F\) ;
    \item pour tout \(\lambda\in\K\) et tout \(u\in F\), on a \(\lambda u\in F\).
\end{enumerate}
\end{theoreme}

\begin{proof}
Supposons d’abord que \(F\) soit un sous-espace vectoriel. Alors il est non vide car il contient son vecteur nul, et il est stable par addition et multiplication par un scalaire.

Réciproquement, supposons \(F\neq\varnothing\), stable par addition et par multiplication par un scalaire.

Comme \(F\neq\varnothing\), il existe \(u\in F\). Alors
\[
0u=0_E\in F,
\]
donc \(F\) contient le vecteur nul de \(E\).

Soit maintenant \(u\in F\). Comme \(-1\in\K\), on a
\[
(-1)u=-u\in F.
\]
Donc tout vecteur de \(F\) admet son opposé dans \(F\).

Les autres propriétés nécessaires sont héritées de \(E\). Ainsi \(F\) est un espace vectoriel pour les lois induites, donc un sous-espace vectoriel de \(E\).
\end{proof}

\begin{corollaire}
Une partie \(F\subset E\) est un sous-espace vectoriel de \(E\) si et seulement si :
\begin{itemize}
    \item \(0_E\in F\) ;
    \item pour tous \(u,v\in F\) et tous \(\lambda,\mu\in\K\),
    \[
    \lambda u+\mu v\in F.
    \]
\end{itemize}
\end{corollaire}

\begin{proof}
Si \(F\) est un sous-espace vectoriel, alors il contient \(0_E\), est stable par multiplication par un scalaire et par addition ; donc \(\lambda u+\mu v\in F\).

Réciproquement, si \(0_E\in F\) et si \(\lambda u+\mu v\in F\) pour tous \(u,v\in F\), alors :
\begin{itemize}
    \item en prenant \(\lambda=\mu=1\), on obtient la stabilité par addition ;
    \item en prenant \(v=0_E\) et \(\mu=0\), on obtient la stabilité par multiplication par un scalaire.
\end{itemize}
Donc \(F\) est un sous-espace vectoriel.
\end{proof}

\subsection{Exemples de sous-espaces vectoriels}

\begin{exemple}
Dans \(\R^2\), l’ensemble
\[
F=\{(x,y)\in\R^2\mid y=2x\}
\]
est un sous-espace vectoriel.

En effet :
\begin{itemize}
    \item \((0,0)\in F\) car \(0=2\times 0\) ;
    \item si \((x,2x)\in F\) et \((x',2x')\in F\), alors
    \[
    (x,2x)+(x',2x')=(x+x',2(x+x'))\in F ;
    \]
    \item pour tout \(\lambda\in\R\),
    \[
    \lambda(x,2x)=(\lambda x,2\lambda x)\in F.
    \]
\end{itemize}
Donc \(F\) est un sous-espace vectoriel.
\end{exemple}

\begin{exemple}
Dans \(\R^3\), l’ensemble
\[
P=\{(x,y,z)\in\R^3\mid x+y-z=0\}
\]
est un sous-espace vectoriel.
\end{exemple}

\begin{exemple}
Dans \(\mathcal{F}(\R,\R)\), l’ensemble des fonctions paires est un sous-espace vectoriel.

En effet :
\begin{itemize}
    \item la fonction nulle est paire ;
    \item la somme de deux fonctions paires est paire ;
    \item le produit d’une fonction paire par un réel est encore paire.
\end{itemize}
\end{exemple}

\subsection{Exemples de parties qui ne sont pas des sous-espaces}

\begin{exemple}
Dans \(\R^2\), l’ensemble
\[
A=\{(x,y)\in\R^2\mid y=2x+1\}
\]
n’est pas un sous-espace vectoriel car
\[
(0,0)\notin A.
\]
Plus conceptuellement, cette droite est affine, mais ne passe pas par l’origine.
\end{exemple}

\begin{exemple}
Dans \(\R\), l’intervalle \([0,+\infty[\) n’est pas un sous-espace vectoriel car il n’est pas stable par multiplication par \(-1\).
\end{exemple}

\section{Intersection et somme de sous-espaces}

\subsection{Intersection de sous-espaces}

\begin{proposition}
L’intersection de deux sous-espaces vectoriels de \(E\) est un sous-espace vectoriel de \(E\).
\end{proposition}

\begin{proof}
Soient \(F\) et \(G\) deux sous-espaces vectoriels de \(E\).  
Comme \(0_E\in F\) et \(0_E\in G\), on a
\[
0_E\in F\cap G.
\]
Donc \(F\cap G\neq \varnothing\).

Soient \(u,v\in F\cap G\). Alors \(u,v\in F\) et \(u,v\in G\). Comme \(F\) et \(G\) sont des sous-espaces,
\[
u+v\in F \quad \text{et} \quad u+v\in G,
\]
donc
\[
u+v\in F\cap G.
\]

Soit \(\lambda\in\K\) et \(u\in F\cap G\). Alors \(u\in F\) et \(u\in G\), donc
\[
\lambda u\in F \quad \text{et} \quad \lambda u\in G.
\]
Ainsi
\[
\lambda u\in F\cap G.
\]
Donc \(F\cap G\) est un sous-espace vectoriel.
\end{proof}

\begin{remarque}
Cette propriété reste vraie pour l’intersection d’une famille quelconque de sous-espaces, à condition que l’intersection soit prise dans un même espace vectoriel.
\end{remarque}

\subsection{Union de sous-espaces}

\begin{remarque}
L’union de deux sous-espaces n’est en général pas un sous-espace vectoriel.
\end{remarque}

\begin{exemple}
Dans \(\R^2\), posons
\[
F=\{(x,0)\mid x\in\R\},
\qquad
G=\{(0,y)\mid y\in\R\}.
\]
Ce sont deux sous-espaces vectoriels. Mais leur union n’est pas un sous-espace, car
\[
(1,0)\in F\subset F\cup G,
\qquad
(0,1)\in G\subset F\cup G,
\]
mais
\[
(1,0)+(0,1)=(1,1)\notin F\cup G.
\]
\end{exemple}

\subsection{Somme de deux sous-espaces}

\begin{definition}
Soient \(F\) et \(G\) deux sous-espaces vectoriels de \(E\). On appelle \textbf{somme} de \(F\) et \(G\) l’ensemble
\[
F+G=\{u+v\mid u\in F,\ v\in G\}.
\]
\end{definition}

\begin{proposition}
La somme \(F+G\) de deux sous-espaces vectoriels de \(E\) est un sous-espace vectoriel de \(E\).
\end{proposition}

\begin{proof}
Comme \(0_E\in F\) et \(0_E\in G\),
\[
0_E=0_E+0_E\in F+G.
\]

Soient \(x,y\in F+G\). Il existe \(u,u'\in F\) et \(v,v'\in G\) tels que
\[
x=u+v,\qquad y=u'+v'.
\]
Alors
\[
x+y=(u+u')+(v+v').
\]
Comme \(F\) et \(G\) sont des sous-espaces,
\[
u+u'\in F,\qquad v+v'\in G.
\]
Donc
\[
x+y\in F+G.
\]

Soit \(\lambda\in\K\) et \(x=u+v\in F+G\). Alors
\[
\lambda x=\lambda(u+v)=\lambda u+\lambda v.
\]
Or \(\lambda u\in F\) et \(\lambda v\in G\), donc
\[
\lambda x\in F+G.
\]
Ainsi \(F+G\) est un sous-espace vectoriel.
\end{proof}

\begin{exemple}
Dans \(\R^2\), si
\[
F=\{(x,0)\mid x\in\R\}
\quad \text{et} \quad
G=\{(0,y)\mid y\in\R\},
\]
alors
\[
F+G=\R^2.
\]
\end{exemple}

\section{Combinaisons linéaires}

\subsection{Définition}

\begin{definition}
Soient \(u_1,\dots,u_p\in E\).  
Une \textbf{combinaison linéaire} de \(u_1,\dots,u_p\) est tout vecteur de la forme
\[
\lambda_1u_1+\cdots+\lambda_pu_p,
\qquad
\text{où } \lambda_1,\dots,\lambda_p\in\K.
\]
\end{definition}

\begin{remarque}
Le mot \og linéaire \fg{} signifie ici qu’on ne fait intervenir que des additions et des multiplications scalaires, pas de produits entre vecteurs ni d’exposants sur les vecteurs.
\end{remarque}

\begin{exemple}
Dans \(\R^3\), le vecteur
\[
(1,2,3)
\]
est une combinaison linéaire de
\[
(1,0,0),\ (0,1,0),\ (0,0,1),
\]
car
\[
(1,2,3)=1(1,0,0)+2(0,1,0)+3(0,0,1).
\]
\end{exemple}

\subsection{Sous-espace engendré}

\begin{definition}
Soient \(u_1,\dots,u_p\in E\).  
Le \textbf{sous-espace vectoriel engendré} par \(u_1,\dots,u_p\), noté
\[
\mathrm{Vect}(u_1,\dots,u_p),
\]
est l’ensemble de toutes les combinaisons linéaires de \(u_1,\dots,u_p\) :
\[
\mathrm{Vect}(u_1,\dots,u_p)=
\left\{
\lambda_1u_1+\cdots+\lambda_pu_p
\mid
\lambda_1,\dots,\lambda_p\in\K
\right\}.
\]
\end{definition}

\begin{proposition}
\(\mathrm{Vect}(u_1,\dots,u_p)\) est un sous-espace vectoriel de \(E\).
\end{proposition}

\begin{proof}
On a
\[
0_E=0u_1+\cdots+0u_p\in \mathrm{Vect}(u_1,\dots,u_p).
\]
Donc cet ensemble est non vide.

Soient
\[
x=\lambda_1u_1+\cdots+\lambda_pu_p
\quad \text{et} \quad
y=\mu_1u_1+\cdots+\mu_pu_p
\]
dans \(\mathrm{Vect}(u_1,\dots,u_p)\). Alors
\[
x+y=(\lambda_1+\mu_1)u_1+\cdots+(\lambda_p+\mu_p)u_p,
\]
donc \(x+y\in \mathrm{Vect}(u_1,\dots,u_p)\).

Pour \(\alpha\in\K\),
\[
\alpha x=(\alpha\lambda_1)u_1+\cdots+(\alpha\lambda_p)u_p,
\]
donc \(\alpha x\in \mathrm{Vect}(u_1,\dots,u_p)\).

Ainsi \(\mathrm{Vect}(u_1,\dots,u_p)\) est un sous-espace vectoriel.
\end{proof}

\begin{remarque}
\(\mathrm{Vect}(u_1,\dots,u_p)\) est le plus petit sous-espace vectoriel contenant \(u_1,\dots,u_p\).
\end{remarque}

\begin{proposition}
Soit \(F\) un sous-espace vectoriel de \(E\).  
Alors
\[
u_1,\dots,u_p\in F
\quad \Longrightarrow \quad
\mathrm{Vect}(u_1,\dots,u_p)\subset F.
\]
\end{proposition}

\begin{proof}
Comme \(F\) est stable par multiplication par un scalaire puis par addition, toute combinaison linéaire de \(u_1,\dots,u_p\) appartient à \(F\).
\end{proof}

\subsection{Exemples de sous-espaces engendrés}

\begin{exemple}
Dans \(\R^2\),
\[
\mathrm{Vect}((1,2))
=
\{(\lambda,2\lambda)\mid \lambda\in\R\}.
\]
C’est la droite vectorielle de direction \((1,2)\).
\end{exemple}

\begin{exemple}
Dans \(\R^3\),
\[
\mathrm{Vect}((1,0,0),(0,1,0))
=
\{(x,y,0)\mid x,y\in\R\}.
\]
C’est le plan vectoriel \(z=0\).
\end{exemple}

\section{Familles finies de vecteurs et premières méthodes}

\subsection{Comment montrer qu’un ensemble est un sous-espace ?}

\begin{methode}
Pour montrer qu’une partie \(F\subset E\) est un sous-espace vectoriel, on suit presque toujours ce schéma :
\begin{enumerate}
    \item vérifier que \(0_E\in F\) ;
    \item prendre deux éléments quelconques \(u,v\in F\) ;
    \item montrer que pour tous \(\lambda,\mu\in\K\),
    \[
    \lambda u+\mu v\in F.
    \]
\end{enumerate}
\end{methode}

\begin{exemple}
Montrons que
\[
F=\{(x,y,z)\in\R^3\mid x-2y+z=0\}
\]
est un sous-espace vectoriel de \(\R^3\).

\textbf{Étape 1 :} \((0,0,0)\in F\), car
\[
0-2\cdot 0+0=0.
\]

\textbf{Étape 2 :} soient \(u=(x,y,z)\in F\) et \(v=(x',y',z')\in F\), et soient \(\lambda,\mu\in\R\).

Comme \(u,v\in F\),
\[
x-2y+z=0
\qquad \text{et} \qquad
x'-2y'+z'=0.
\]

Considérons
\[
\lambda u+\mu v=(\lambda x+\mu x',\, \lambda y+\mu y',\, \lambda z+\mu z').
\]
Alors
\[
(\lambda x+\mu x')-2(\lambda y+\mu y')+(\lambda z+\mu z')
=
\lambda(x-2y+z)+\mu(x'-2y'+z')=0.
\]
Donc \(\lambda u+\mu v\in F\), et \(F\) est un sous-espace vectoriel.
\end{exemple}

\subsection{Comment montrer qu’un ensemble n’est pas un sous-espace ?}

\begin{methode}
Pour montrer qu’une partie \(A\subset E\) n’est pas un sous-espace vectoriel, il suffit de montrer qu’une des propriétés nécessaires échoue :
\begin{itemize}
    \item \(0_E\notin A\) ;
    \item absence de stabilité par addition ;
    \item absence de stabilité par multiplication par un scalaire.
\end{itemize}
\end{methode}

\begin{exemple}
Dans \(\R^2\), l’ensemble
\[
A=\{(x,y)\in\R^2\mid xy=1\}
\]
n’est pas un sous-espace vectoriel, car
\[
(0,0)\notin A.
\]
\end{exemple}

\begin{exemple}
Dans \(\R^2\), l’ensemble
\[
B=\{(x,y)\in\R^2\mid x\geq 0\}
\]
n’est pas un sous-espace vectoriel, car il n’est pas stable par multiplication par \(-1\).
\end{exemple}

\section{Analyse-synthèse dans les espaces vectoriels}

\subsection{Principe}

\begin{remarque}
La méthode d’analyse-synthèse est très utile en algèbre linéaire, notamment pour :
\begin{itemize}
    \item décrire un sous-espace engendré ;
    \item caractériser l’ensemble des solutions d’un problème ;
    \item trouver une écriture d’un vecteur sous forme de combinaison linéaire.
\end{itemize}
\end{remarque}

\begin{exemple}
Décrivons
\[
\mathrm{Vect}\big((1,1),(1,-1)\big)\subset \R^2.
\]

\textbf{Analyse :} un élément de cet ensemble est de la forme
\[
\lambda(1,1)+\mu(1,-1)=(\lambda+\mu,\lambda-\mu).
\]
Si on note
\[
x=\lambda+\mu,\qquad y=\lambda-\mu,
\]
on peut espérer obtenir tous les couples \((x,y)\in\R^2\).

\textbf{Synthèse :} soit \((x,y)\in\R^2\). Cherchons \(\lambda,\mu\in\R\) tels que
\[
(\lambda+\mu,\lambda-\mu)=(x,y).
\]
Cela revient à résoudre
\[
\lambda+\mu=x,\qquad \lambda-\mu=y.
\]
En ajoutant et en soustrayant, on obtient
\[
\lambda=\frac{x+y}{2},\qquad \mu=\frac{x-y}{2}.
\]
Donc tout \((x,y)\in\R^2\) s’écrit comme combinaison linéaire de \((1,1)\) et \((1,-1)\).

Ainsi
\[
\mathrm{Vect}\big((1,1),(1,-1)\big)=\R^2.
\]
\end{exemple}

\section{Erreurs classiques}

\begin{itemize}
    \item Oublier qu’un sous-espace vectoriel doit contenir le vecteur nul.
    \item Confondre une droite affine et une droite vectorielle.
    \item Penser que l’union de deux sous-espaces est toujours un sous-espace.
    \item Montrer seulement la stabilité par addition sans vérifier la stabilité par multiplication par un scalaire.
    \item Écrire \(\mathrm{Vect}(u,v)\) sans préciser qu’il s’agit de l’ensemble de toutes les combinaisons linéaires de \(u\) et \(v\).
    \item Confondre vecteurs et scalaires.
\end{itemize}

\section{Résumé du chapitre}

\begin{itemize}
    \item Un espace vectoriel est un ensemble muni d’une addition et d’une multiplication par un scalaire, vérifiant des axiomes précis.
    \item Les espaces fondamentaux sont \(\K^n\), \(M_{n,p}(\K)\), \(\K[X]\) et les espaces de fonctions.
    \item Un sous-espace vectoriel est une partie stable par combinaison linéaire et contenant le vecteur nul.
    \item L’intersection de sous-espaces est un sous-espace.
    \item L’union de sous-espaces n’est en général pas un sous-espace.
    \item La somme \(F+G\) de deux sous-espaces est un sous-espace.
    \item Le sous-espace engendré par des vecteurs est l’ensemble de leurs combinaisons linéaires.
\end{itemize}

\section*{Exercices d’application immédiate}

\begin{enumerate}[label=\textbf{Exercice \arabic*.}]
    \item Montrer que
    \[
    F=\{(x,y,z)\in\R^3\mid x+y+z=0\}
    \]
    est un sous-espace vectoriel de \(\R^3\).

    \item Montrer que
    \[
    A=\{(x,y)\in\R^2\mid x+y=1\}
    \]
    n’est pas un sous-espace vectoriel de \(\R^2\).

    \item Déterminer explicitement
    \[
    \mathrm{Vect}((1,0),(0,1)) \subset \R^2.
    \]

    \item Déterminer explicitement
    \[
    \mathrm{Vect}((1,2,0),(0,1,3)) \subset \R^3.
    \]

    \item Montrer que l’ensemble des fonctions impaires de \(\mathcal{F}(\R,\R)\) est un sous-espace vectoriel.

    \item Soient
    \[
    F=\{(x,y,z)\in\R^3\mid x+y=0\},
    \qquad
    G=\{(x,y,z)\in\R^3\mid z=0\}.
    \]
    Décrire \(F\cap G\) et \(F+G\).
\end{enumerate}

\end{document}